% status: FULL
\chapter{The Feed-Forward Network and the Residual Stream}\label{ch:feed-forward}
\begin{ChapterIntent}
Attention moves information between positions. A feed-forward network
changes the features at one position. We will build that second operation
with small matrices, see why the nonlinear step matters, and connect it
to the vector that carries information through the whole decoder.
\end{ChapterIntent}

\section{Start with one token vector}
Take an illustrative token representation $\mathbf x=[1,-1]$. Its two
coordinates are features, not probabilities. We want a transformation that
can react differently to positive and negative combinations of those
features. Start by expanding two coordinates into three:
\[
 \mathbf h=[x_0,\ x_1,\ x_0-x_1]=[1,-1,2].
\]
This is a linear projection. It has a matrix representation, but the
coordinate expressions already tell us every multiplication. Apply ReLU,
which keeps positive values and replaces negative values by zero:
\[
 \operatorname{ReLU}(\mathbf h)=[1,0,2].
\]
Then combine the first two activated coordinates into one output and the
last two into another. Writing $\mathbf a=\operatorname{ReLU}(\mathbf h)$,
\[
 \Delta\mathbf x=[a_0+a_1,\ a_1+a_2]=[1,2].
\]
The output width is again two. The wider intermediate representation has
done its job and can be discarded after the output is calculated.

A dense feed-forward network applies the same learned matrices separately
to every token row. It does not select another token or read its value
vector. If one row changes, this operator alone cannot change the outputs
of the other rows. Attention in a later layer can spread that change.
Keeping the two roles separate is the easiest way to read a decoder diagram.

\section{Why not use a single matrix?}
Two linear maps without an intervening nonlinearity can be combined into
one linear map. If row-vector activations are the convention, then
$(\mathbf x\mathbf W_1)\mathbf W_2=\mathbf x(\mathbf W_1\mathbf W_2)$.
Changing the intermediate width would not change this fact. The activation
prevents us from collapsing the general transformation into one fixed
matrix.

Try changing the input to $[-1,1]$. Before activation the hidden vector
becomes $[-1,1,-2]$; after ReLU it is $[0,1,0]$. The update becomes
$[1,1]$, not the negative of $[1,2]$. A linear transformation would have
obeyed $f(-\mathbf x)=-f(\mathbf x)$. This small counterexample shows
exactly where additional expressive power enters.

For a sequence matrix $\mathbf X\in\mathbb R^{S\times d}$, a conventional
two-matrix feed-forward sublayer is
\[
 \operatorname{FFN}(\mathbf X)=
 \phi(\mathbf X\mathbf W_{\rm up}+\mathbf b_{\rm up})
 \mathbf W_{\rm down}+\mathbf b_{\rm down}.
\]
The up matrix has shape $d\times d_{\rm ff}$ and the down matrix
$d_{\rm ff}\times d$. Biases broadcast over token rows. The activation
$\phi$ acts element by element. These shapes are part of the definition;
the word \emph{linear} in a library API may include a bias.

\section{Follow the residual, not only the update}
The feed-forward output usually does not replace the incoming vector.
It is added to a \emph{residual stream}: the model-width vector passed
from one sublayer to the next. In our first example,
\[
 \mathbf x_{\rm next}=\mathbf x+\Delta\mathbf x
                    =[1,-1]+[1,2]=[2,1].
\]
The input, update and result are three different objects. Overwriting the
input with a normalized copy before saving the residual can quietly change
the model even though every array has the right shape.

In the pre-normalized decoder used in this book, the two additions are
\[
 \mathbf r=\mathbf x+
   \operatorname{Attention}(\operatorname{Norm}(\mathbf x)),\qquad
 \mathbf x_{\rm next}=\mathbf r+
   \operatorname{FFN}(\operatorname{Norm}(\mathbf r)).
\]
The attention expression includes its Q/K/V and output projections.
The second normalization reads $\mathbf r$, not the original $\mathbf x$.
Each sublayer reads a normalized view of the current residual but adds
its update to the unnormalized residual.

Other architectures put normalization elsewhere. That is not a cosmetic
difference: changing the order changes the numerical function. We choose
one convention explicitly so that the implementation has one answer to
test. A checkpoint loader must select the convention required by its model.

\EngineFigure{foundation-residual-ffn}{The pre-normalized residual block
built here. The vertical path carries the model-width representation.
Each side branch computes an update; a plus sign adds it to the unchanged
input of that sublayer. Attention mixes positions; the FFN mixes features
within each position.}{fig:foundation-residual-ffn}

\section{Build a gate with two projections}
Many decoder models use a gated feed-forward network instead of the
two-matrix ReLU construction. It has two input projections and one output
projection. One branch supplies features; the other modulates them.
For the bias-free SwiGLU variant \cite{shazeer2020glu},
\[
 \mathbf g=\operatorname{SiLU}(\mathbf x\mathbf W_g),\quad
 \mathbf u=\mathbf x\mathbf W_u,\quad
 \Delta\mathbf x=(\mathbf g\odot\mathbf u)\mathbf W_d,
 \qquad \operatorname{SiLU}(z)=\frac{z}{1+e^{-z}}.
\]
The symbol $\odot$ means coordinatewise multiplication, not a dot
product. Both intermediate branches have width $d_{\rm ff}$. The down
projection returns to width $d$, making residual addition possible.
The gate is not a probability. SiLU can be negative, and its positive
output is not bounded by one. It changes the strength and sign of
features; it does not make a yes/no choice of which token to attend to.
Nor is this a mixture-of-experts router. Every hidden coordinate in this
dense FFN is computed.

Use the same input $[1,-1]$ with projections chosen to produce
gate preactivations $[0,1]$ and up values $[2,3]$. SiLU gives
$[0,1/(1+e^{-1})]$. Multiplying branches yields
$[0,3/(1+e^{-1})]\approx[0,2.193176]$.
Let the down projection produce $[h_0-h_1,\ 0.5h_0+h_1]$. The update is
$[-2.193176,2.193176]$ and, if added directly to our original illustrative
input, the result is approximately $[-1.193176,1.193176]$.
This arithmetic omits normalization to isolate the FFN; the complete
decoder will normalize its FFN input.

\section{Count weights and work before optimizing}
Without biases, an ordinary two-matrix FFN has $2dd_{\rm ff}$ weights.
A three-matrix gated FFN has $3dd_{\rm ff}$ weights. With $d=4$ and
$d_{\rm ff}=6$, the gated version has $3\cdot4\cdot6=72$ weights.
At four bytes each, its matrix payload is 288 bytes.

Using the convention that a multiply--add counts as two floating-point
operations, those projections perform about $6dd_{\rm ff}$ FLOPs per
token, or 144 in the example. That estimate excludes the activation,
elementwise product, normalization and residual addition. It is a useful
matrix-work count, not the exact instruction count.

A sequence of $S$ tokens applies the same matrices to $S$ rows.
The tokens do not communicate within the FFN, so an implementation can
batch those rows into a matrix product. The parameter count stays fixed
while activation storage and arithmetic grow with the number of rows.
This is one of the first places where model structure creates an engine
optimization: the same weight can serve several token rows.

The hidden width need not equal four times model width. Architectures
choose it along with the activation and projection count. Comparing a
gated FFN and a two-matrix FFN at the same hidden width therefore compares
different parameter budgets. Read the model configuration rather than
inferring one width from a familiar architecture name.

\section{Implement a scalar reference}
Run the example:
\begin{verbatim}
python3 code/foundations/lesson.py ffn
\end{verbatim}
It prints the ReLU hidden row, its down-projected update and the SwiGLU
update. The returned vectors are updates, not residual sums. Add the
input by hand to check that distinction.

The \texttt{swiglu} function in \texttt{code/foundations/model.py}
performs two calls to \texttt{linear}, activates the gate, multiplies
matching coordinates and performs the third projection. Weight matrices
in this program store \emph{output rows}: \texttt{weights[j][i]}
multiplies input coordinate $i$ into output coordinate $j$. This is the
transpose of the row-vector mathematical storage convention above.
The function's scalar definition, not the printed orientation of an
array, determines whether a transpose is needed.

SiLU also needs a numerically sensible sigmoid evaluation. For a large
negative input, directly evaluating $e^{-z}$ can overflow. The teaching
function uses $z e^z/(1+e^z)$ for negative $z$, the algebraically
equivalent form whose exponential is bounded. The test suite includes
$z=-1000$ and requires a finite result.

\section{Test a transformation, not a name}
A useful FFN test suite asks several separate questions. Does the up
projection produce the expected coordinates? Does the activation handle
negative, zero and positive inputs? Do gate and up branches have equal
width? Does the down projection return model width? Is the residual
added exactly once?

Our hand examples answer the first three numerical questions with fixed
expected values. The full decoder tests answer the composition question
by comparing execution paths. In a production kernel, compare optimized
outputs to a scalar reference under an explicit tolerance. Merely
observing that both implementations are named SwiGLU establishes nothing.

One especially useful deliberate break is to add the original input
twice. The output remains finite and shape-correct. The small residual
example immediately exposes it: $[1,-1]+[1,2]=[2,1]$, whereas the double
addition gives $[3,0]$. Another break is swapping the gate and up
projections. Since SiLU is nonlinear, the two branches are not generally
interchangeable.

\section{Make a branch disappear on purpose}
There is a useful boundary test for the bias-free gated network.
Set every gate-projection weight to zero. Every gate preactivation
becomes zero, SiLU of zero is zero, and the coordinatewise product
is zero regardless of the up branch. The down projection therefore
returns a zero update. With residual addition, the incoming residual
passes through unchanged.

This test checks more than a convenient special case. A nonzero
result indicates an unexpected bias, a stale intermediate buffer,
a wrong branch or an extra operation. The expected answer does not
depend on a second implementation of SwiGLU. It follows directly
from zero multiplied by any finite value.

Now restore the gate and set the down matrix to zero. Again the
update must vanish. In a complete layer, this removes the FFN
update but leaves the earlier attention update intact. It is a
controlled way to isolate one sublayer while debugging composition,
not a proposal to remove learned FFNs from a real model.

\section{Where this goes next}
We now have the two sublayers of a decoder block and their residual
connections. Attention brings context into each position. The FFN
transforms the resulting features. Normalization supplies the scale
convention, while the residual carries the representation forward.
The next chapter stacks these blocks, projects the final representation
to vocabulary logits, and runs the first complete generation loop.
Nothing about that assembly requires a large model or a GPU.
\section{Worked problems}
\subsection*{1. Reverse the input}
\textbf{Problem.} Apply the chapter's ReLU FFN to $[-1,1]$.
\textbf{Solution.} The expanded row is $[-1,1,-2]$, activation gives
$[0,1,0]$, the update is $[1,1]$ and the residual result is $[0,2]$.
\subsection*{2. Price the projections}
\textbf{Problem.} Count bias-free SwiGLU weights and matrix FLOPs for
$d=4$, $d_{\rm ff}=6$, and three input tokens.
\textbf{Solution.} There are 72 weights and 288 F32 bytes.
Matrix work is approximately $3\cdot6\cdot4\cdot6=432$ FLOPs.
Weights are shared, so the three-token batch does not triple parameter storage.
\subsection*{3. A local operator}
\textbf{Problem.} How can you test that an FFN does not mix token positions?
\textbf{Solution sketch.} Run two rows, change only one row, and rerun.
The untouched row's output must be unchanged. Also compare the batched
operation with separate row calls under the intended numerical tolerance.
\Needspace{8\baselineskip}
\subsection*{4. Which vector receives the update?}
\textbf{Problem.} For the normalized trace, incoming residual is $[3,4]$
and attention update is $[1,-2]$. Find the first residual and the FFN
input under unit RMS gains and $\epsilon=0.5$.
\textbf{Solution.} Add to the unnormalized residual to get $[4,2]$.
Its mean square is ten, giving FFN input $[4,2]/\sqrt{10.5}$.
Using the original normalized attention input as the residual changes
the model's function.
\Needspace{8\baselineskip}
\subsection*{5. Equal matrix budgets}
\textbf{Problem.} A bias-free two-matrix FFN has model width six and hidden
width twelve. Find a gated hidden width with the same matrix parameter
count and conventional matrix FLOPs per token.
\textbf{Solution.} The original has $2\cdot6\cdot12=144$ parameters.
Solve $3\cdot6\cdot f=144$ to get $f=8$. Both have 288 conventional
matrix FLOPs per token. Nonlinear work and runtime need not be equal.
\Needspace{8\baselineskip}
\subsection*{6. Retire the hidden tail}
\textbf{Problem.} A five-wide hidden axis is streamed in tiles of two.
How many products are in each tile, and what must remain live?
\textbf{Solution.} Tile lengths are two, two and one. The maximum product
list length is two. All model-width output accumulators remain live
until every hidden product contributes. Omitting the one-entry tail
changes the update; it is not harmless padding.

